% ============================================================
% Embedded bibliography: this .tex file can generate Bibliography.bib automatically.
% Figure PDFs remain separate submission assets under Figures/.
% ============================================================
\begin{filecontents*}{Bibliography.bib}
@article{Makridakis2022M5,
  author  = {Makridakis, Spyros and Spiliotis, Evangelos and Assimakopoulos, Vassilios},
  title   = {The M5 accuracy competition: Results, findings, and conclusions},
  journal = {International Journal of Forecasting},
  year    = {2022},
  volume  = {38},
  number  = {4},
  pages   = {1346--1364},
  doi     = {10.1016/j.ijforecast.2021.11.013}
}

@article{Theodorou2025Inventory,
  author  = {Theodorou, Evangelos and Spiliotis, Evangelos and Assimakopoulos, Vassilios},
  title   = {Forecast accuracy and inventory performance: Insights on their relationship from the M5 competition data},
  journal = {European Journal of Operational Research},
  year    = {2025},
  volume  = {322},
  number  = {2},
  pages   = {414--426},
  doi     = {10.1016/j.ejor.2024.12.033}
}

@article{Breiman2001RF,
  author  = {Breiman, Leo},
  title   = {Random forests},
  journal = {Machine Learning},
  year    = {2001},
  volume  = {45},
  number  = {1},
  pages   = {5--32},
  doi     = {10.1023/A:1010933404324}
}

@inproceedings{Chen2016XGBoost,
  author    = {Chen, Tianqi and Guestrin, Carlos},
  title     = {XGBoost: A scalable tree boosting system},
  booktitle = {Proceedings of the 22nd ACM SIGKDD International Conference on Knowledge Discovery and Data Mining},
  year      = {2016},
  pages     = {785--794},
  doi       = {10.1145/2939672.2939785}
}

@article{Nsir2025XGBoost,
  author  = {Nsir, Noreddin and Alzubi, Ahmad Bassam and Adegboye, Oluwatayomi Rereloluwa},
  title   = {Enhancing sustainable supply chain performance prediction using an augmented algorithm-optimized XGBoost in Industry 4.0 contexts},
  journal = {Sustainability},
  year    = {2025},
  volume  = {17},
  number  = {22},
  pages   = {10344},
  doi     = {10.3390/su172210344}
}

@inproceedings{Deepak2025RFLinear,
  author    = {Deepak, R. and Sathyanarayanan, R. and Arunnehru, J.},
  title     = {Demand and sales forecasting using random forest and linear regression},
  booktitle = {Proceedings of the International Conference on Data Science and Applications},
  year      = {2025},
  publisher = {Springer Nature},
  pages     = {525--540},
  doi       = {10.1007/978-981-96-1981-8\_42}
}

@article{Theodoridis2025LSTMixer,
  author  = {Theodoridis, Georgios and Tsadiras, Athanasios},
  title   = {Retail demand forecasting: A comparative analysis of deep neural networks and the proposal of LSTMixer, a linear model extension},
  journal = {Information},
  year    = {2025},
  volume  = {16},
  number  = {7},
  pages   = {596},
  doi     = {10.3390/info16070596}
}

@article{Yang2025MARL,
  author  = {Yang, Yongbin and Wang, Mengdie and Wang, Jiyuan and Li, Pan and Zhou, Mengjie},
  title   = {Multi-agent deep reinforcement learning for integrated demand forecasting and inventory optimization in sensor-enabled retail supply chains},
  journal = {Sensors},
  year    = {2025},
  volume  = {25},
  number  = {8},
  pages   = {2428},
  doi     = {10.3390/s25082428}
}

@article{Pensa2025ExplainXGB,
  author  = {Pensa, Ruggero G. and Crombach, Anton and Peignier, Sergio and Rigotti, Christophe},
  title   = {Explaining random forest and XGBoost with shallow decision trees by co-clustering feature importance},
  journal = {Machine Learning},
  year    = {2025},
  volume  = {114},
  number  = {12},
  doi     = {10.1007/s10994-025-06932-9}
}

@article{Long2025ProbRetail,
  author  = {Long, Xueying and Bui, Quang and Oktavian, Grady and Schmidt, Daniel F. and Bergmeir, Christoph and Godahewa, Rakshitha and Lee, Seong Per and Zhao, Kaifeng and Condylis, Paul},
  title   = {Scalable probabilistic forecasting in retail with gradient boosted trees: A practitioner's approach},
  journal = {International Journal of Production Economics},
  year    = {2025},
  volume  = {279},
  pages   = {109449},
  doi     = {10.1016/j.ijpe.2024.109449}
}

@article{Liu2025Stockout,
  author  = {Liu, Yang and Kalaitzi, Dimitra and Wang, Michael and Papanagnou, Christos},
  title   = {A machine learning approach to inventory stockout prediction},
  journal = {Journal of Digital Economy},
  year    = {2025},
  volume  = {4},
  pages   = {144--155},
  doi     = {10.1016/j.jdec.2025.06.002}
}

@article{Samal2025Ensemble,
  author  = {Samal, Tejaswini and Ghosh, Anupam},
  title   = {Ensemble-based predictive analytics for demand forecasting in multi-channel retailing},
  journal = {Expert Systems with Applications},
  year    = {2025},
  volume  = {299},
  pages   = {130212},
  doi     = {10.1016/j.eswa.2025.130212}
}

@article{Fildes2022RetailForecasting,
  author  = {Fildes, Robert and Ma, Shaohui and Kolassa, Stephan},
  title   = {Retail forecasting: Research and practice},
  journal = {International Journal of Forecasting},
  year    = {2022},
  volume  = {38},
  number  = {4},
  pages   = {1283--1318},
  doi     = {10.1016/j.ijforecast.2019.06.004}
}

@inproceedings{Gonzalez2025AccuracyPerformance,
  author    = {Gonzalez, Pablo and Rodriguez, Pedro and Luna, Marcos and Ponte, Borja},
  title     = {From accuracy to performance: Advanced demand forecasting in supply chains},
  booktitle = {Proceedings of the International Conference on Industrial Engineering and Operations Management},
  year      = {2025},
  publisher = {Springer Nature},
  pages     = {407--412},
  doi       = {10.1007/978-3-031-82334-3\_71}
}

@article{Abolghasemi2025LLMForecasting,
  author  = {Abolghasemi, Mahdi and Ganbold, Odkhishig and Rotaru, Kristian},
  title   = {Humans vs. large language models: Judgmental forecasting in an era of advanced AI},
  journal = {International Journal of Forecasting},
  year    = {2025},
  volume  = {41},
  number  = {2},
  pages   = {631--648},
  doi     = {10.1016/j.ijforecast.2024.07.003}
}

@article{Giannopoulos2025Intermittent,
  author  = {Giannopoulos, Panagiotis G. and Dasaklis, Thomas K. and Tsantilis, Ioannis and Patsakis, Constantinos},
  title   = {Machine learning algorithms in intermittent demand forecasting: A review},
  journal = {International Journal of Production Research},
  year    = {2025},
  doi     = {10.1080/00207543.2025.2578701}
}

@article{Hu2025AutoMLRetail,
  author  = {Hu, Hengzhi and Tan, Dan and Thaichon, Park and Wang, Bolin and Zhu, Zhicheng},
  title   = {Grid-based market sales forecasting for retail businesses using automated machine learning and geospatial intelligence},
  journal = {Expert Systems with Applications},
  year    = {2025},
  volume  = {284},
  pages   = {127869},
  doi     = {10.1016/j.eswa.2025.127869}
}

@article{Petropoulos2025Occam,
  author  = {Petropoulos, Fotios and Grushka-Cockayne, Yael and Siemsen, Enno and Spiliotis, Evangelos},
  title   = {Wielding Occam's razor: Fast and frugal retail forecasting},
  journal = {Journal of the Operational Research Society},
  year    = {2025},
  volume  = {76},
  number  = {8},
  pages   = {1564--1583},
  doi     = {10.1080/01605682.2024.2421339}
}

@article{Punia2026HumanML,
  author  = {Punia, Sushil},
  title   = {Medium- to long-term demand forecasting in retail and manufacturing organizations: Integration of machine learning, human judgment, and interval variable},
  journal = {Journal of Forecasting},
  year    = {2026},
  volume  = {45},
  number  = {1},
  pages   = {122--134},
  doi     = {10.1002/for.70030}
}

@article{Mohammed2026Uncertainty,
  author  = {Mohammed, Zizi and Chafi, Anas and El Hammoume, Mohammed},
  title   = {A hybrid learning framework for forecasting uncertainty and adaptive inventory planning in retail supply chains},
  journal = {Supply Chain Analytics},
  year    = {2026},
  volume  = {13},
  pages   = {100180},
  doi     = {10.1016/j.sca.2025.100180}
}

@article{Sattar2025SCM,
  author  = {Sattar, M. U. and Dattana, V. and Hasan, R. and Mahmood, S. and Khan, H. W. and Hussain, S.},
  title   = {Enhancing supply chain management: A comparative study of machine learning models for forecasting and ESG performance},
  journal = {Sustainability},
  year    = {2025},
  volume  = {17},
  number  = {13},
  pages   = {5772},
  doi     = {10.3390/su17135772}
}

@article{Ermochenko2026Ecommerce,
  author  = {Ermochenko, S. A.},
  title   = {An integrated approach to demand forecasting and inventory optimization in e-commerce},
  journal = {Technoeconomics},
  year    = {2026},
  volume  = {5},
  number  = {1},
  pages   = {75--84},
  doi     = {10.57809/2026.5.1.16.7}
}

@article{Feda2026SHAPRetail,
  author  = {Feda, Afi Kekeli and Caglar, Ersin},
  title   = {A metaheuristic-optimized extreme learning machine with SHAP-driven interpretability for retail demand forecasting},
  journal = {Discover Computing},
  year    = {2026},
  volume  = {29},
  pages   = {257},
  doi     = {10.1007/s10791-026-10075-3}
}

@article{Shahzad2025ThreeSTF,
  author  = {Shahzad, Muhammad Ahmad and Fong, Titus Hei Yeung},
  title   = {Demand forecasting for retail using three-S temporal fusion (3STF) network},
  journal = {Strategic Business Research},
  year    = {2025},
  volume  = {2},
  number  = {1},
  pages   = {100030},
  doi     = {10.1016/j.sbr.2025.100030}
}

@article{Chintapanti2025PrGB,
  author  = {Chintapanti, Anusha and Maiti, Sandipan},
  title   = {Optimizing retail inventory and sales through advanced time series forecasting using fine tuned PrGB regressor},
  journal = {IEEE Access},
  year    = {2025},
  volume  = {13},
  pages   = {165328--165341},
  doi     = {10.1109/ACCESS.2025.3605229}
}

@article{Ahmed2026RetailARIMA,
  author  = {Ahmed, Ali F. and Mohammed, Osamah S. and Khan, Nokhaiz Tariq and Al Hanbali, Ahmad and Mohammed, Awsan},
  title   = {Improving forecast accuracy for retail supply chains using ARIMA and machine learning approaches},
  journal = {Transportation Research Procedia},
  year    = {2026},
  volume  = {96},
  pages   = {524--531},
  doi     = {10.1016/j.trpro.2026.04.024}
}

@article{Petropoulos2022ForecastingTheory,
  author  = {Petropoulos, Fotios and Apiletti, Daniele and Assimakopoulos, Vassilios and Babai, Mohamed Zied and Barrow, Devon K. and Ben Taieb, Souhaib and Bergmeir, Christoph and Bessa, Ricardo J. and Bijak, Jakub and Boylan, John E. and others},
  title   = {Forecasting: Theory and practice},
  journal = {International Journal of Forecasting},
  year    = {2022},
  volume  = {38},
  number  = {3},
  pages   = {705--871},
  doi     = {10.1016/j.ijforecast.2021.11.001}
}

@article{Lim2021TFT,
  author  = {Lim, Bryan and Arik, Sercan O. and Loeff, Nicolas and Pfister, Tomas},
  title   = {Temporal fusion transformers for interpretable multi-horizon time series forecasting},
  journal = {International Journal of Forecasting},
  year    = {2021},
  volume  = {37},
  number  = {4},
  pages   = {1748--1764},
  doi     = {10.1016/j.ijforecast.2021.03.012}
}

@article{Makridakis2018M4,
  author  = {Makridakis, Spyros and Spiliotis, Evangelos and Assimakopoulos, Vassilios},
  title   = {The M4 competition: Results, findings, conclusion and way forward},
  journal = {International Journal of Forecasting},
  year    = {2018},
  volume  = {34},
  number  = {4},
  pages   = {802--808},
  doi     = {10.1016/j.ijforecast.2018.06.001}
}

@misc{DataCo2019,
  author       = {Constante, Fabian and Silva, Fernando and Pereira, Antonio},
  title        = {DataCo SMART SUPPLY CHAIN FOR BIG DATA ANALYSIS},
  year         = {2019},
  publisher    = {Mendeley Data},
  version      = {5},
  doi          = {10.17632/8gx2fvg2k6.5},
  url          = {https://doi.org/10.17632/8gx2fvg2k6.5}
}

@article{Bergsma2025InventoryML,
 author={Bergsma, Ritsaart and de Ruijt, Corné and Bhulai, Sandjai}, title={A systematic review of machine learning approaches in inventory control optimization}, journal={Operations Research Perspectives}, year={2025}, volume={15}, pages={100367}, doi={10.1016/j.orp.2025.100367}}

@article{Nguyen2025TSMixerInventory,
 author={Nguyen, Thi Ngoc Anh and Nguyen, Thi Xuan Hoa and Tran, Ngoc Thang and Nguyen, Thi Ha and Nguyen, Phuong Anh and Vu, Hai Anh}, title={Optimizing supply chain operations using advanced Time-Series Mixer models for demand forecasting and inventory under uncertain demand}, journal={Expert Systems with Applications}, year={2026}, volume={296}, pages={128955}, doi={10.1016/j.eswa.2025.128955}}

@article{Retraining2025Retail,
 author={Zanotti, Marco}, title={On the retraining frequency of global models in retail demand forecasting}, journal={Machine Learning with Applications}, year={2025}, volume={22}, pages={100769}, doi={10.1016/j.mlwa.2025.100769}}

@article{Nguyen2025SupplyChainML,
 author={Nguyen, Thi Ngoc Anh and others}, title={Enhancing sustainable supply chain forecasting using machine learning for sales prediction}, journal={Procedia Computer Science}, year={2025}, volume={252}, pages={470--479}, doi={10.1016/j.procs.2025.01.006}}

@article{Federated2025Demand,
 author={Qi, Hang and Luo, Jieping and Li, Qiyue and Wu, Jingjin}, title={A comparative trade-off analysis on accuracy and efficiency for federated learning in demand forecasting}, journal={Applied Soft Computing}, year={2025}, volume={182}, pages={113561}, doi={10.1016/j.asoc.2025.113561}}

@article{AnticipatoryShipping2025,
 author={Ren, Xinxin and Zeng, Kuan}, title={Data-driven analysis on inventory problem for anticipatory shipping}, journal={Computers \& Industrial Engineering}, year={2025}, volume={203}, pages={111038}, doi={10.1016/j.cie.2025.111038}}

@article{RIFT2026SupplyChain,
 author={Wang, Yanchun}, title={Construction of supply chain demand forecasting algorithm based on time series data analysis and improvement of its management efficiency}, journal={Procedia Computer Science}, year={2026}, volume={281}, pages={484--493}, doi={10.1016/j.procs.2026.04.240}}

@article{Naser2025SPINEX,
 author={Naser, Ahmad Z. and Naser, M. Z.},
 title={SPINEX-TimeSeries: Similarity-based predictions with explainable neighbors exploration for time series and forecasting problems},
 journal={Computers \& Industrial Engineering},
 year={2025},
 volume={200},
 pages={110812},
 doi={10.1016/j.cie.2024.110812}}

@article{OptimizedHyperparams2025,
 author={Vhatkar, Manjunath S. and Mahajan, Pramod Sanjay and Raut, Rakesh D. and Cheikhrouhou, Naoufel and Ghoshal, Subhash}, title={Optimized hyperparameters for retail sales forecasting using grid search}, journal={Engineering Applications of Artificial Intelligence}, year={2025}, volume={158}, pages={111472}, doi={10.1016/j.engappai.2025.111472}}

@article{Surana2026SparseTransformer,
 author={Surana, Ishan and Ahuja, Divej and Singh, Udai Pratap and Kaliraj, S. and Reddy, G. Pradeep},
 title={An autoregressive lag-aware sparse transformer framework for demand forecasting},
 journal={Discover Artificial Intelligence}, year={2026}, volume={6}, pages={465},
 doi={10.1007/s44163-026-01204-4}}
\end{filecontents*}

% ============================================================
% C&IE submission manuscript generated from the thesis-aligned manuscript draft.
% Numerical findings are preserved; only manuscript presentation and figure paths are consolidated.
% Graphical abstract is intentionally not embedded here: Elsevier requests it as a separate submission file.
% Computers & Industrial Engineering - Journal Manuscript
% Based on the final UE Master Thesis:
% "A Data-Driven Framework for Retail Demand Forecasting
%  and Financial Impact Evaluation Using the M5 Benchmark Dataset"
% ============================================================

\documentclass[preprint,12pt]{elsarticle}

\usepackage{amssymb}
\usepackage{amsmath}
\usepackage{booktabs}
\usepackage{array}
\usepackage{multirow}
\usepackage{adjustbox}
\usepackage{graphicx}
\usepackage{xcolor}
\usepackage{url}
\usepackage{placeins}

\journal{Computers \& Industrial Engineering}

% Safe figure command: the manuscript compiles even if a figure has
% not yet been uploaded. Replace the filename with the final PDF name.
\newcommand{\safeincludegraphics}[2][]{%
  \IfFileExists{\detokenize{#2}}{\includegraphics[#1]{#2}}{%
    \fbox{\parbox[c][0.16\textheight][c]{0.92\linewidth}{%
      \centering Figure file not yet uploaded\\[2mm]
      \texttt{\detokenize{#2}}}}}}

\begin{document}

\begin{frontmatter}

\title{A Data-Driven Framework for Retail Demand Forecasting and Financial and Inventory Scenario Analysis Using the M5 Benchmark Dataset}

\author[label1]{Mohamed Meraj Mansur Sharif\corref{cor1}}\ead{mohamed.sharif@ue-germany.de}
\author[label1]{Raja Hashim Ali}
\author[label1]{Talha Ali Khan}
\cortext[cor1]{Corresponding author.}

\affiliation[label1]{organization={Department of Business, University of Europe for Applied Sciences},
            addressline={Think Campus, Konrad-Zuse-Ring 11},
            city={Potsdam},
            postcode={14469},
            state={Brandenburg},
            country={Germany}}

\begin{abstract}
Retail demand forecasting is important for inventory planning, financial analysis, and supply-chain decision-making, but predictive accuracy alone does not show how forecast errors translate into operational consequences.
This study develops a reproducible machine-learning framework that connects retail demand forecasting with financial and inventory scenario analysis using the M5 Forecasting dataset.
The study selects the 500 highest-volume item--store series using training-period demand, integrates sales, calendar, and weekly price data, and engineers lag, rolling-window, pricing, and calendar-derived features.
Linear Regression, Random Forest, and XGBoost are evaluated under identical conditions using a chronological 28-day hold-out containing 14,000 observations, while LightGBM and ARIMA are used as supporting analyses.
The full feature configuration improves Random Forest MAE from 5.842 to 5.608, while the 7-day rolling mean accounts for 79.9\% of Random Forest feature importance.
Among the three core models, XGBoost achieves the strongest row-level performance with an MAE of 5.582, RMSE of 8.690, and $R^2$ of 0.705.
Using LightGBM forecasts for downstream decision analysis, the illustrative scenario produces approximately \$335,984.96 in fulfilled-demand revenue, \$6,020.06 in stockout-cost exposure, and \$3,273.91 in holding-cost exposure.
Estimated safety stock increases from approximately 12.0 units at a 70\% service target to 37.8 units at 95\%.
External validation on the DataCo SMART SUPPLY CHAIN dataset produces an MAE of 50.32 and RMSE of 61.25 on its own scale.
The framework therefore provides a transparent link between predictive performance, uncertainty, inventory requirements, and financial consequences while explicitly distinguishing empirical results from illustrative decision scenarios.
\end{abstract}

\begin{highlights}
\item A multi-source framework links retail forecasts to inventory decisions.
\item XGBoost achieves MAE 5.582, RMSE 8.690 and $R^2=0.705$ on the M5 hold-out.
\item Recent demand history provides the dominant predictive signal in selected series.
\item Safety-stock scenarios rise from 12.0 units at 70\% service to 37.8 units at 95\%.
\item External DataCo validation supports cautious transferability of the feature strategy.
\end{highlights}

\begin{keyword}
Demand forecasting
\sep Machine learning
\sep Retail supply chain
\sep Inventory optimisation
\sep Financial impact analysis
\sep M5 forecasting dataset
\end{keyword}

\end{frontmatter}


% ============================================================
\section{Introduction}
% ============================================================

Retail demand forecasting sits near the centre of procurement, replenishment, workforce planning, and service-level decisions in modern supply chains \cite{Fildes2022RetailForecasting,Petropoulos2022ForecastingTheory}.
Forecast errors have asymmetric operational consequences because under-forecasting can contribute to stockouts and lost sales, whereas over-forecasting can leave capital tied up in excess inventory and increase holding costs.
The value of a forecasting system therefore depends not only on predictive accuracy but also on how its errors interact with inventory and cost decisions.
Retailers increasingly have access to transactional, pricing, calendar, and product-level information that can be incorporated into machine-learning forecasting pipelines \cite{Nguyen2025SupplyChainML,Makridakis2022M5}.
Tree-based ensemble methods are particularly attractive for structured retail data because they can model nonlinear relationships while accommodating heterogeneous engineered features \cite{Hu2025AutoMLRetail,Samal2025Ensemble}.
At the same time, the forecasting literature continues to debate the relative value of statistical, machine-learning, and deep-learning approaches across datasets and forecast horizons \cite{Makridakis2018M4,Lim2021TFT,Petropoulos2025Occam}.
The present study focuses on a transparent tabular machine-learning setting rather than attempting to claim universal superiority for any single model family.
The empirical analysis uses the M5 Forecasting benchmark and evaluates the predictive and downstream implications of forecasts within one consistent experimental pipeline.

\subsection{Gap Analysis}

A recurring limitation in the retail forecasting literature is the separation between statistical forecast accuracy and the operational or financial value of those forecasts.
Studies frequently report MAE, RMSE, or related accuracy measures without translating prediction errors into inventory exposure, revenue implications, or service-level requirements \cite{Theodorou2025Inventory,Gonzalez2025AccuracyPerformance}.
Conversely, studies that connect forecasts with inventory decisions often use different datasets, different forecasting procedures, or simulation settings that make direct comparison difficult.
Recent work has strengthened probabilistic forecasting and forecast-driven inventory analysis, but the two evaluation layers are still commonly treated as separate methodological exercises \cite{Long2025ProbRetail,Mohammed2026Uncertainty}.
This separation matters because a small improvement in RMSE does not necessarily produce a proportional reduction in stockout or holding costs.
The operational consequence of an error depends on demand level, forecast uncertainty, selling price, cost assumptions, lead time, and the inventory policy used to translate the prediction into action.
There is therefore a need for transparent empirical frameworks in which model accuracy, uncertainty, and decision-oriented quantities are derived from the same held-out observations.
This study addresses that gap by carrying a single M5-based pipeline from feature engineering and forecasting through financial and inventory scenario analysis.

\subsection{Research Questions}

The study investigates one main research question and six supporting questions:

\begin{enumerate}
    \item \textbf{RQ1:} How accurately can machine-learning models forecast retail product demand using a multi-source feature set built from the M5 dataset, and how much of that accuracy comes from combining historical, pricing, and calendar information rather than from any single source alone?
    \item \textbf{RQ2:} Does multi-source feature engineering improve forecasting performance relative to a narrower feature set, and which engineered feature families contribute the most predictive signal?
    \item \textbf{RQ3:} Which of Linear Regression, Random Forest, and XGBoost provides the most reliable predictive performance under identical experimental conditions?
    \item \textbf{RQ4:} How does forecast output translate into inventory-related financial performance, including revenue, stockout cost, and holding cost?
    \item \textbf{RQ5:} Can demand forecasting support inventory optimisation by quantifying the trade-off between target service level and safety-stock investment?
    \item \textbf{RQ6:} How does forecasting uncertainty relate to stockout risk, and what does that relationship imply for safety-stock policy?
    \item \textbf{RQ7:} Does forecasting performance remain stable across low-, medium-, and high-demand regimes, and does the underlying feature-engineering approach generalise beyond the M5 benchmark?
\end{enumerate}

\subsection{Problem Statement}

The problem addressed in this study is that retail organisations need forecasting systems that are evaluated not only on statistical accuracy but also on the financial and operational consequences of that accuracy.
The study addresses this problem for the M5 Forecasting dataset by engineering a multi-source feature set from historical sales, calendar, and pricing data.
Three forecasting algorithms of increasing complexity are trained under identical inputs and a chronological validation design.
The resulting predictions are then translated into revenue, stockout cost, holding cost, and service-level-dependent safety-stock requirements.
Forecast uncertainty is further examined by separating observations into uncertainty tiers and relating error magnitude to stockout risk.
Demand-regime analysis is used to identify whether forecast performance is uniform across low-, medium-, and high-demand series.
An external DataCo experiment provides a separate robustness check without treating the different target scale as directly comparable with M5.
The resulting framework therefore answers both how accurately demand can be forecast and what that forecast accuracy means for inventory-oriented decisions.

\subsection{Novelty of This Study}

The contribution is methodological rather than a claim of state-of-the-art leaderboard performance.
The study combines several evaluation layers that are commonly reported separately and keeps the underlying observations and validation horizon fixed.

\begin{itemize}
    \item \textbf{Unified evaluation:} Linear Regression, Random Forest, and XGBoost are compared using the same preprocessing, feature matrix, chronological split, and test observations.
    \item \textbf{Decision translation:} The same held-out forecasting outputs are translated into revenue, stockout-cost, holding-cost, and safety-stock scenarios.
    \item \textbf{Uncertainty-aware analysis:} Forecast-error magnitude is used to differentiate safety-stock requirements and identify higher-risk observations.
    \item \textbf{Regime analysis:} Forecast accuracy is evaluated separately across low-, medium-, and high-demand regimes.
    \item \textbf{External validation:} The core lag/rolling feature-engineering strategy is examined on the DataCo SMART SUPPLY CHAIN dataset while explicitly avoiding direct metric equivalence across different target scales.
\end{itemize}

\subsection{Significance of Our Work}

The study provides a compact decision-support perspective on retail forecasting in which prediction quality and downstream consequences are evaluated together without conflating them.
The results show that recent demand history carries most of the predictive signal in the selected high-volume M5 series.
The strongest core model is XGBoost, but the performance difference from Random Forest is small.
The downstream analysis shows that forecast uncertainty is unevenly distributed and that safety-stock requirements increase substantially as the target service level becomes more demanding.
The financial quantities are deliberately treated as scenario outputs because retailer-specific procurement, holding, shortage, and lost-sales costs are unavailable in the public benchmark.
This distinction prevents the dollar values from being interpreted as observed Walmart financial performance.
The external validation further indicates that the basic temporal feature-engineering strategy can be transferred to a structurally different supply-chain dataset, although the experiment is not a controlled cross-dataset benchmark.
Together, these findings support a transparent forecasting framework in which model accuracy is considered as one component of a broader inventory decision problem.

% ============================================================
\section{Literature Review}
% ============================================================

Retail demand forecasting spans statistical time-series models, tree-based machine learning, deep-learning sequence models, feature-engineering approaches, and decision-oriented inventory analytics.
Classical statistical methods such as ARIMA remain attractive because of their transparency and low computational requirements, while machine-learning models can incorporate external variables such as price and calendar information \cite{Petropoulos2022ForecastingTheory,Makridakis2018M4}.
Tree-based ensembles have performed strongly on structured retail datasets, including the M5 competition \cite{Breiman2001RF,Chen2016XGBoost,Makridakis2022M5}.
Deep-learning models such as LSTM and Transformer architectures provide alternative mechanisms for learning temporal dependencies and probabilistic representations \cite{Theodoridis2025LSTMixer,Lim2021TFT}.
A parallel literature emphasises feature engineering and the value of multi-source predictors, particularly lag and rolling-window statistics \cite{Nguyen2025SupplyChainML,Bergsma2025InventoryML}.
Finally, inventory research increasingly examines how forecasts can be translated into stocking decisions, with recent work emphasising the non-linear relationship between forecast accuracy and inventory performance \cite{Bergsma2025InventoryML,Theodorou2025Inventory}.
The following subsections position the present study within these strands, and Table~\ref{tab:literature} summarises the comparative positioning.

\subsection{Tree-Based Ensemble Learning}

Random Forest uses bootstrap aggregation and random feature selection to reduce the variance of decision-tree predictors \cite{Breiman2001RF}.
Its suitability for tabular forecasting arises from its ability to represent nonlinear interactions without requiring extensive functional assumptions.
XGBoost extends the tree-ensemble family through sequential gradient boosting and regularisation \cite{Chen2016XGBoost}.
The M5 competition demonstrated the effectiveness of carefully engineered boosted-tree systems at large retail forecasting scale \cite{Makridakis2022M5}.
More recent retail studies have continued to report strong performance from simplified or tuned tree-based methods, suggesting that feature engineering and implementation discipline can matter as much as model complexity \cite{OptimizedHyperparams2025,Deepak2025RFLinear,Sattar2025SCM}.
The present study therefore compares Linear Regression, Random Forest, and XGBoost directly rather than assuming that the most complex model must perform best.
LightGBM is included only as a supporting model for diagnostics and downstream scenario analysis.
This separation keeps the headline comparison interpretable while allowing a closely related boosting model to support the decision-oriented experiments.

\subsection{Deep Learning and Hybrid Forecasting}

LSTM networks model sequential dependencies through recurrent memory mechanisms and have been widely applied to demand forecasting \cite{Theodoridis2025LSTMixer}.
Transformer architectures use attention mechanisms to represent long-range dependencies and have become increasingly prominent in time-series forecasting \cite{Lim2021TFT}.
DeepAR provides a probabilistic recurrent alternative that learns across multiple related time series and produces predictive distributions rather than only point forecasts \cite{Long2025ProbRetail}.
Recent retail studies have also explored Temporal Fusion Transformers and other attention-based architectures on M5 and related datasets \cite{Theodoridis2025LSTMixer,Shahzad2025ThreeSTF,Yang2025MARL}.
Recent C\&IE work also shows the value of explicitly combining multiple time scales before forecasting, reinforcing the broader importance of temporal feature representation in supply-chain demand prediction \cite{Naser2025SPINEX}.
These approaches can be powerful when data volume, sequence structure, and tuning resources justify their complexity \cite{Yang2025MARL,Shahzad2025ThreeSTF,Feda2026SHAPRetail}.
However, their additional computational and modelling requirements can make them less attractive for smaller post-filtering samples or applications where transparent tabular feature effects are important.
The present study consequently does not claim that tree-based models are universally superior to deep learning.
Instead, it evaluates a deliberately bounded model set suited to the available M5 analytical sample and the study's emphasis on reproducibility and decision translation.

\subsection{Feature Engineering and Multi-Source Data Integration}

Feature engineering transforms raw transactional and temporal information into variables that expose useful predictive structure to a learning algorithm.
Lagged demand and rolling-window statistics are particularly important because they encode recent demand history while avoiding the need for a model to infer every temporal relationship from raw dates.
Reviews of supply-chain forecasting literature identify feature engineering and the integration of temporal, pricing, and promotional information as recurring determinants of forecasting quality \cite{Nguyen2025SupplyChainML,Bergsma2025InventoryML}.
Systematic mapping work also notes that incomplete reporting of feature construction can make otherwise strong forecasting studies difficult to reproduce \cite{Bergsma2025InventoryML}.
Feature attribution methods such as SHAP provide more formal explanations for tree-based predictions, while simpler impurity-based feature importance can provide a descriptive diagnostic \cite{Pensa2025ExplainXGB}.
Recent work also compares forecasting approaches under distributed learning constraints using Walmart sales data, illustrating that accuracy and computational efficiency can be considered jointly \cite{Federated2025Demand}.
The present study uses Random Forest impurity-based importance because the objective is to identify the relative contribution of engineered predictors rather than to provide a full local explanation framework.
The results are interpreted cautiously because correlated lag and rolling features can share predictive information.
This feature-centric design also provides a direct test of whether adding price and calendar variables improves a narrower historical-demand baseline.

\subsection{Forecast-Driven Inventory Optimisation}

The translation of forecasts into inventory decisions is rooted in classical operations research.
Safety stock provides a buffer against demand and forecast uncertainty, while service level represents the desired probability or proportion of demand satisfied without stockout \cite{Bergsma2025InventoryML}.
The newsvendor problem formalises the trade-off between ordering too little and ordering too much under uncertain demand \cite{Bergsma2025InventoryML}.
Machine-learning research has increasingly revisited this problem by using predictive variables to estimate demand or order quantities directly \cite{Bergsma2025InventoryML}.
Recent work has also embedded inventory cost objectives more directly into forecasting models, reinforcing the value of evaluating forecast accuracy alongside inventory consequences \cite{Nguyen2025TSMixerInventory,AnticipatoryShipping2025}.
Recent empirical work using M5 has shown that forecast accuracy and inventory performance are related but not necessarily one-to-one \cite{Theodorou2025Inventory}.
Probabilistic forecasting research similarly argues that uncertainty information can be important for downstream inventory planning \cite{Long2025ProbRetail,Mohammed2026Uncertainty,Retraining2025Retail,Bergsma2025InventoryML}.
The present study retains a transparent two-step forecast-then-translate approach rather than directly optimising the inventory decision.
This design keeps the statistical accuracy measures and scenario calculations separately interpretable while still allowing their relationship to be examined.

\subsection{Hierarchical Structure and Robustness}

Retail demand data are naturally hierarchical, with products aggregated through categories, stores, and geographical levels.
Hierarchical forecasting research has developed reconciliation methods that seek coherent forecasts across aggregation levels \cite{RIFT2026SupplyChain}.
The present study does not perform hierarchical reconciliation and therefore does not claim coherence across the full M5 hierarchy.
Instead, it focuses on 500 high-volume item--store series and examines robustness through demand-regime analysis and external validation.
This narrower design is appropriate for the research questions because it isolates the relationship between engineered features, forecast accuracy, and downstream decision scenarios.
The limitation is that the results cannot automatically be transferred to slow-moving or intermittent-demand series.
Research on intermittent demand and seasonal retail forecasting reinforces the importance of considering demand regime and data sparsity when evaluating forecasting systems \cite{Giannopoulos2025Intermittent,RIFT2026SupplyChain}.
The robustness analyses in this study are therefore presented as empirical diagnostics rather than as proof of universal generalisation \cite{Nguyen2025SupplyChainML,Ermochenko2026Ecommerce}.
Recent lag-aware transformer research further illustrates the continuing expansion of demand-forecasting methods beyond the bounded tabular model set used here \cite{Surana2026SparseTransformer}.

\begin{table*}[!t]
\centering
\caption{Comparative positioning of the present study relative to representative demand-forecasting and inventory-analytics research. The comparison is conceptual because the cited studies use different datasets, targets, horizons, and evaluation protocols.}
\label{tab:literature}
\scriptsize
\begin{adjustbox}{width=\textwidth}
\begin{tabular}{p{2.5cm} p{2.7cm} p{3.1cm} p{3.0cm} p{4.2cm}}
\toprule
\textbf{Study} & \textbf{Focus} & \textbf{Data / Method} & \textbf{Main contribution} & \textbf{Difference from present study} \\
\midrule
Vhatkar et al. \cite{OptimizedHyperparams2025} & Retail forecasting & Tree-based methods with explanatory variables & Shows that simplified tree methods can be highly competitive & Does not provide the present study's financial and safety-stock translation \\
Nsir et al. \cite{Nsir2025XGBoost} & Supply-chain prediction & Optimised XGBoost & Demonstrates high predictive accuracy after optimisation & Focuses on prediction rather than an integrated inventory-cost analysis \\
Bergsma et al. \cite{Bergsma2025InventoryML} & Literature review & ML/DL supply-chain forecasting & Synthesises the role of multi-source features & Review-level contribution without a new empirical pipeline \\
Theodorou et al. \cite{Theodorou2025Inventory} & Forecast--inventory relationship & M5 competition data & Shows that forecast accuracy and inventory performance are related non-linearly & Does not provide the same head-to-head three-model and uncertainty workflow \\
Long et al. \cite{Long2025ProbRetail} & Probabilistic retail forecasting & Gradient-boosted trees & Connects uncertainty-aware forecasting with retail applications & Uses a probabilistic framework rather than the present post-hoc error-based safety-stock analysis \\
Liu et al. \cite{Liu2025Stockout} & Stockout prediction & Machine learning & Identifies near-term forecasting information as relevant to stockout prediction & Focuses on stockout prediction rather than revenue, holding cost, and safety stock together \\
\textbf{Present study} & \textbf{Forecasting + decision translation} & \textbf{M5 top-500 series; LR, RF, XGBoost; supporting LightGBM/ARIMA} & \textbf{Unifies accuracy, financial, inventory, uncertainty, regime, and external-validation analyses} & \textbf{Single benchmark and bounded modelling sample; scenario costs and lead time are illustrative} \\
\bottomrule
\end{tabular}
\end{adjustbox}
\end{table*}

% ============================================================
\section{Methodology}
% ============================================================

\subsection{Detailed Methodology}

The research design combines data selection, multi-source integration, feature engineering, model training, chronological validation, and downstream financial and inventory analysis within a single executable notebook.
The notebook was executed in a Kaggle environment using Python and standard open-source data-science libraries.
The three core models are kept separate from the supporting LightGBM and ARIMA analyses.
All headline M5 models use the same training observations, feature matrix, and 28-day test horizon.
The downstream decision analyses use the same M5 test observations and LightGBM predictions rather than introducing a second test sample.
This design ensures that the accuracy and decision-oriented layers remain connected while their purposes remain distinct.
The workflow is shown in Figure~\ref{fig:workflow}, and the multi-source forecasting framework is shown in Figure~\ref{fig:framework}.
The external DataCo experiment is treated as a robustness check rather than as a directly comparable benchmark \cite{DataCo2019}.

\begin{figure*}[!t]
\centering
\includegraphics[height=0.82\textheight,keepaspectratio]{Figures/workflow.pdf}
\caption{Overall research workflow used in the study, showing M5 series selection using training-period demand, data integration, preprocessing, lag and rolling feature engineering, chronological model training and validation, financial and inventory scenario analysis, uncertainty and demand-regime diagnostics, and external validation using the DataCo SMART SUPPLY CHAIN dataset.}
\label{fig:workflow}
\end{figure*}

\begin{figure*}[!t]
\centering
\includegraphics[width=0.95\textwidth]{Figures/multimodal_framework.pdf}
\caption{Multi-source forecasting framework linking heterogeneous retail inputs, feature engineering, model evaluation, and decision-oriented financial and inventory analysis.}
\label{fig:framework}
\end{figure*}


\subsection{Dataset}

The primary empirical source is the M5 Forecasting Accuracy competition dataset \cite{Makridakis2022M5}.
The dataset contains daily Walmart unit sales for 3,049 products across ten stores in California, Texas, and Wisconsin, represented as 30,490 item--store time series.
The raw sales table contains 30,490 rows and 1,919 columns, while the calendar and selling-price tables contain 1,969 rows by 14 columns and 6,841,121 rows by 4 columns, respectively.
The analysis selects the 500 highest-volume item--store series using only training-period sales.
This prevents test-period demand from influencing series selection.
The selected series are reshaped into 956,500 product-day observations before calendar and price information are merged.
The final model-ready feature matrix contains ten numerical predictors. Table~\ref{tab:rawdata} summarises the primary data sources used before preprocessing.

An independent DataCo SMART SUPPLY CHAIN dataset is used for external validation, where transactions are aggregated to daily order counts and a separate chronological 80/20 evaluation is performed \cite{DataCo2019}.

\begin{table}[!t]
\centering
\caption{Raw M5 data files used before preprocessing and model-series selection.}
\label{tab:rawdata}
\begin{tabular}{lrr}
\toprule
\textbf{File} & \textbf{Rows} & \textbf{Columns} \\
\midrule
Sales (`sales\_train\_validation.csv') & 30,490 & 1,919 \\
Calendar (`calendar.csv') & 1,969 & 14 \\
Prices (`sell\_prices.csv') & 6,841,121 & 4 \\
\bottomrule
\end{tabular}
\end{table}

\subsection{Data Preprocessing}

Date fields are converted to datetime format and decomposed into year, month, ISO week, day of week, and quarter.
For each item--store series, 7-day and 28-day lag variables and 7-day and 28-day trailing means are constructed.
The rolling means are shifted by one day before the rolling operation so that the current target is not included in its own predictors.
Selling prices are forward-filled and backward-filled within each item--store series before the period-over-period price-change feature is calculated.
Rows are not removed merely because event-name or event-type information is unavailable in the final model matrix.
Instead, only rows lacking required numerical modelling variables are removed.
The resulting structural loss occurs at the beginning of each series because lag and rolling variables require historical observations.
The initial 956,500 product-day observations are reduced to 942,500 analytical rows.

\subsection{Feature Engineering}

The final feature matrix contains ten numerical predictors grouped into three information families.
Historical-demand features are \texttt{lag\_7}, \texttt{lag\_28}, \texttt{rolling\_mean\_7}, and \texttt{rolling\_mean\_28}.
Pricing features are \texttt{sell\_price} and \texttt{price\_change}.
Calendar-derived features are \texttt{month}, \texttt{week}, \texttt{dayofweek}, and \texttt{quarter}.
Product and store identifiers are retained for grouping and aggregation but are not supplied directly as predictors.
RQ1 evaluates progressively richer feature configurations.
RQ2 uses Random Forest impurity-based feature importance to identify the relative contribution of the engineered variables.
Because several lag and rolling features are correlated, importance values are interpreted as descriptive rather than causal measures of information contribution.

\subsection{Forecasting Models}

Linear Regression provides the simplest baseline and models demand as a weighted combination of the engineered predictors.
Random Forest uses 50 trees, maximum depth 16, minimum leaf size 2, a 50\% bootstrap fraction, random seed 42, and two parallel jobs.
XGBoost uses 100 estimators, learning rate 0.1, maximum depth 6, random seed 42, and a squared-error regression objective.
LightGBM is used as a supporting model with 100 estimators and random seed 42 for M5 downstream analyses.
ARIMA(5,1,0) is used only as a supporting statistical benchmark on the aggregate daily series.
The three core models are trained on the same ten-feature matrix and evaluated on exactly the same test observations.
LightGBM and ARIMA are therefore not mixed into the headline three-model ranking.
This separation is maintained throughout the Results and Discussion sections.

\subsection{Experimental Settings}

The experimental settings were fixed before the final chronological evaluation so that the three core models were compared under identical data and test conditions.
Random Forest, XGBoost, and LightGBM used the configurations reported in Table~\ref{tab:models}.
The supporting ARIMA analysis was fitted only to the aggregate daily series and is not included in the row-level model ranking.

\begin{table}[!t]
\centering
\caption{Model configurations used in the experiments. LightGBM and ARIMA are supporting models and are not part of the headline three-model comparison.}
\label{tab:models}
\begin{tabular}{lccc}
\toprule
\textbf{Parameter} & \textbf{Random Forest} & \textbf{XGBoost} & \textbf{LightGBM} \\
\midrule
Number of estimators & 50 & 100 & 100 \\
Learning rate & -- & 0.1 & Default \\
Maximum tree depth & 16 & 6 & Default \\
Minimum samples per leaf & 2 & -- & -- \\
Bootstrap fraction & 0.5 & -- & -- \\
Random seed & 42 & 42 & 42 \\
\bottomrule
\end{tabular}
\end{table}

\subsection{Chronological Validation}

The final analytical dataset is ordered by date and divided chronologically.
The last 28 calendar days, from 28 March 2016 through 24 April 2016, form the test horizon.
The preceding observations form the training period from 26 February 2011 through 27 March 2016.
This produces 928,500 training rows and 14,000 test rows, corresponding to 500 selected item--store series over 28 test days.
A chronological split is essential because a random row-wise split could allow information from later dates to influence the evaluation of earlier dates.
The 28-day horizon also mirrors the forecasting horizon of the M5 competition.
All three core models are therefore evaluated under the same temporal conditions.

\subsection{Evaluation Metrics}

Forecast accuracy is evaluated using MAE, RMSE, and the coefficient of determination $R^2$.
For observed demand $y_i$ and predicted demand $\hat{y}_i$ over $n$ test observations, MAE is defined as

\begin{equation}
\mathrm{MAE}
=
\frac{1}{n}\sum_{i=1}^{n}|y_i-\hat{y}_i|.
\label{eq:mae}
\end{equation}

RMSE is defined as

\begin{equation}
\mathrm{RMSE}
=
\sqrt{
\frac{1}{n}
\sum_{i=1}^{n}
(y_i-\hat{y}_i)^2
}.
\label{eq:rmse}
\end{equation}

The coefficient of determination is

\begin{equation}
R^2
=
1-
\frac{
\sum_{i=1}^{n}(y_i-\hat{y}_i)^2
}{
\sum_{i=1}^{n}(y_i-\bar{y})^2
}.
\label{eq:r2}
\end{equation}

MAE measures average absolute forecast error in demand units, RMSE gives greater weight to large errors, and $R^2$ describes the proportion of observed demand variance explained by the model \cite{Petropoulos2022ForecastingTheory}.
The three measures are reported together because they capture different characteristics of the error distribution.

\subsection{Financial Translation}

The downstream financial analysis uses LightGBM predictions from the same 14,000 M5 test observations.
Fulfilled-demand revenue is calculated as the minimum of actual and forecast demand multiplied by selling price.
Stockout-cost exposure is calculated as unmet demand multiplied by selling price and a 10\% cost factor.
Holding-cost exposure is calculated as over-forecast demand multiplied by selling price and a 5\% cost factor.
These percentages are illustrative assumptions rather than observed Walmart accounting parameters.
The resulting dollar quantities are therefore scenario outputs rather than estimates of actual Walmart profit, revenue, or operating cost.
This distinction is essential because the public M5 dataset does not contain retailer-specific procurement, shortage, holding, labour, or logistics cost structures.
The financial analysis is intended to demonstrate how forecast errors can be translated into decision-oriented quantities under transparent assumptions.
It does not constitute a complete profit optimisation model.

\subsection{Inventory and Uncertainty Analysis}

Safety stock is estimated using

\begin{equation}
SS=z_{\alpha}\sigma_e\sqrt{L},
\label{eq:safety}
\end{equation}

where $z_{\alpha}$ is the standard-normal quantile associated with the target service level, $\sigma_e$ is the standard deviation of signed forecast error, and $L$ is the assumed lead time.
The main service-level analysis evaluates targets from 70\% to 95\% using a seven-day lead time.
The uncertainty analysis divides the 14,000 observations into three equal-sized tertiles using absolute forecast error.
Tier-specific safety stock is then calculated using each tier's error standard deviation at a 95\% service target and seven-day lead time.
This procedure links inventory buffers directly to observed forecast uncertainty.
It remains a simplified scenario analysis because actual lead-time distributions, replenishment constraints, service-level definitions, and product substitution are not available in M5.

\subsection{External Validation}

External validation uses the DataCo SMART SUPPLY CHAIN FOR BIG DATA ANALYSIS dataset.
The structured data contain 180,519 transaction rows in the notebook.
Transactions are aggregated to daily order counts and four historical-demand features are constructed.
LightGBM is evaluated using a chronological 80/20 split.
The experiment produces MAE 50.32 and RMSE 61.25.
Because the target scale and demand-generating process differ from M5, these metrics are not compared numerically with the M5 metrics.
The purpose of the experiment is instead to determine whether the basic lag and rolling feature-engineering strategy remains usable outside the original benchmark.
The external validation is therefore treated as a robustness check rather than evidence of universal generalisation.

\subsection{Research Ethics and Reproducibility}

The study uses publicly available benchmark and research datasets and does not involve human participants, personally identifiable information, or experimental intervention.
Reproducibility is supported by the single notebook workflow, explicit feature definitions, fixed random seeds for stochastic models, and the fixed chronological hold-out.
The software environment uses Python with pandas, NumPy, scikit-learn, XGBoost, LightGBM, statsmodels, matplotlib, and seaborn.
The model configurations and evaluation design are reported explicitly.
The source notebook and repository should accompany the manuscript according to the journal and institutional submission requirements.
The authors used OpenAI ChatGPT during manuscript preparation for language editing, restructuring, LaTeX formatting, and selected code/notebook debugging. No AI system generated or altered the empirical data, model outputs, research figures, or scientific conclusions; all such outputs were produced from the underlying data and verified by the authors.

% ============================================================
\section{Results}
% ============================================================

All headline M5 results use the same chronological 28-day test horizon containing 14,000 item--store-day observations.
The core comparison consists of Linear Regression, Random Forest, and XGBoost.
LightGBM and ARIMA are reported only as supporting diagnostics.
The downstream financial, inventory, uncertainty, and demand-regime analyses use LightGBM predictions from the same M5 test observations.

\subsection{RQ1: Overall Forecasting Accuracy}

RQ1 evaluates how forecasting accuracy changes as historical, pricing, and calendar information are combined. Table~\ref{tab:rq1} and Figure~\ref{fig:rq1config} report the four feature configurations and their corresponding MAE values.
Four feature configurations are compared using Random Forest: a lag-and-rolling baseline, a price-enhanced configuration, a calendar-enhanced configuration, and the full multi-source feature matrix.
The baseline obtains MAE 5.842.
Adding price information alone changes MAE only slightly to 5.836.
The calendar-enhanced configuration performs worse, with MAE 6.691.
The full multi-source configuration achieves the lowest MAE at 5.608, representing approximately a 4.0\% improvement over the baseline.
These results show that the combined feature matrix performs best in the present experiment, while the magnitude of the improvement is primarily associated with the historical demand features.

\begin{table}[!t]
\centering
\caption{Forecast MAE by feature configuration on the 14,000-observation chronological M5 test set. Lower values indicate lower absolute forecast error.}
\label{tab:rq1}
\begin{tabular}{lc}
\toprule
\textbf{Feature configuration} & \textbf{MAE} \\
\midrule
Baseline: lag and rolling-window features & 5.842 \\
Price-enhanced & 5.836 \\
Calendar-enhanced & 6.691 \\
Full multi-source & \textbf{5.608} \\
\bottomrule
\end{tabular}
\end{table}

\begin{figure}[!t]
\centering
\includegraphics[width=\linewidth]{Figures/rq1_feature_configuration.pdf}
\caption{MAE for the four feature configurations evaluated in RQ1. The full multi-source configuration provides the lowest test-set MAE, while the calendar-enhanced configuration performs worse than the historical-demand baseline.}
\label{fig:rq1config}
\end{figure}
\FloatBarrier

\subsection{RQ2: Contribution of Engineered Features}

RQ2 examines which engineered variables contribute most strongly to the Random Forest predictions. Table~\ref{tab:rq2} and Figure~\ref{fig:rq2importance} report the numerical importance scores and their ranking.
The 7-day rolling mean has an importance score of 0.7991, equivalent to approximately 79.9\% of the reported importance.
The 28-day rolling mean has importance 0.0516.
Day of week, lag 7, lag 28, week, selling price, month, and price change have importance values of 0.0408, 0.0306, 0.0289, 0.0208, 0.0202, 0.0069, and 0.0010, respectively.
The two rolling-window variables together account for approximately 85.5\% of the importance diagnostic.
The correlation analysis shows strong relationships among lag and rolling-window predictors.
This means that the importance values should not be interpreted as independent causal contributions.
Nevertheless, the ranking indicates that recent demand history provides the dominant predictive signal in the selected feature matrix.

\begin{table}[!t]
\centering
\caption{Random Forest impurity-based feature-importance scores for the variables included in the RQ2 diagnostic. The values are descriptive importance measures and should not be interpreted as causal effects.}
\label{tab:rq2}
\begin{tabular}{lc}
\toprule
\textbf{Feature} & \textbf{Importance} \\
\midrule
rolling\_mean\_7 & 0.7991 \\
rolling\_mean\_28 & 0.0516 \\
dayofweek & 0.0408 \\
lag\_7 & 0.0306 \\
lag\_28 & 0.0289 \\
week & 0.0208 \\
sell\_price & 0.0202 \\
month & 0.0069 \\
price\_change & 0.0010 \\
\bottomrule
\end{tabular}
\end{table}

\begin{figure}[!t]
\centering
\includegraphics[width=\linewidth]{Figures/rq2_feature_importance.pdf}
\caption{Random Forest feature-importance ranking for the engineered predictors. The 7-day rolling mean dominates the diagnostic, followed by the 28-day rolling mean, while individual pricing and calendar variables contribute smaller shares.}
\label{fig:rq2importance}
\end{figure}
\FloatBarrier

\subsection{RQ3: Comparative Model Evaluation}

RQ3 compares Linear Regression, Random Forest, and XGBoost under identical training and test conditions. Table~\ref{tab:rq3} and Figure~\ref{fig:rq3comparison} report the core-model results on the identical test observations.
XGBoost achieves the strongest core performance, with MAE 5.582, RMSE 8.690, and $R^2=0.705$.
Random Forest follows closely with MAE 5.608, RMSE 8.753, and $R^2=0.701$.
Linear Regression records MAE 5.935, RMSE 9.141, and $R^2=0.674$.
The supporting LightGBM model achieves MAE 5.584 and RMSE 8.693 on the same row-level M5 observations.
ARIMA achieves MAE 958.511 and RMSE 1,107.919 after the test observations are aggregated to the daily series, so these values are not directly comparable with the item--store row-level metrics.
The core ranking therefore identifies XGBoost as the strongest model among the three primary candidates for this feature matrix and horizon.
The small difference between XGBoost and Random Forest also indicates that the result should not be interpreted as evidence that greater model complexity is universally superior.

\begin{table}[!t]
\centering
\caption{Core model performance on the identical 14,000-observation chronological M5 test set. The supporting LightGBM result is reported separately because it is not part of the headline three-model comparison.}
\label{tab:rq3}
\begin{tabular}{lccc}
\toprule
\textbf{Model} & \textbf{MAE} & \textbf{RMSE} & \textbf{$R^2$} \\
\midrule
Linear Regression & 5.935 & 9.141 & 0.674 \\
Random Forest & 5.608 & 8.753 & 0.701 \\
XGBoost & \textbf{5.582} & \textbf{8.690} & \textbf{0.705} \\
\midrule
Supporting LightGBM & 5.584 & 8.693 & -- \\
\bottomrule
\end{tabular}
\end{table}

\begin{figure}[!t]
\centering
\includegraphics[width=\linewidth]{Figures/rq3_core_model_comparison.pdf}
\caption{Core model comparison on the identical M5 chronological test set. XGBoost provides the lowest MAE and RMSE and the highest $R^2$ among Linear Regression, Random Forest, and XGBoost.}
\label{fig:rq3comparison}
\end{figure}
\FloatBarrier

\subsection{RQ4: Financial Implications of Forecasting}

RQ4 translates the LightGBM forecasts into scenario-level financial quantities using the same 14,000 M5 test observations. Table~\ref{tab:rq4} and Figure~\ref{fig:rq4} report the resulting revenue, stockout-cost, and holding-cost quantities.
Fulfilled-demand revenue is approximately \$335,984.96.
Stockout-cost exposure is approximately \$6,020.06.
Holding-cost exposure is approximately \$3,273.91.
The stockout-cost total is approximately 1.84 times the holding-cost total under the specified assumptions.
This ratio is not a general empirical retail finding because it depends on the realised errors, selling prices, and the assumed 10\% and 5\% cost factors.
The analysis therefore demonstrates the mechanics of financial translation rather than estimating retailer-specific profitability.
Procurement, labour, logistics, and other business costs are outside the calculation.
The quantities should consequently be interpreted as transparent scenarios generated from held-out forecasts.

\begin{table}[!t]
\centering
\caption{Financial implications calculated from LightGBM forecasts over the 14,000-observation M5 test horizon under illustrative unit-cost assumptions.}
\label{tab:rq4}
\begin{tabular}{lr}
\toprule
\textbf{Metric} & \textbf{Value (USD)} \\
\midrule
Fulfilled-demand revenue & 335,984.96 \\
Stockout-cost exposure & 6,020.06 \\
Holding-cost exposure & 3,273.91 \\
\bottomrule
\end{tabular}
\end{table}

\begin{figure}[!t]
\centering
\includegraphics[width=\linewidth]{Figures/rq4_financial_implications.pdf}
\caption{Financial quantities derived from LightGBM forecasts on the M5 test set. Revenue, stockout-cost exposure, and holding-cost exposure are scenario values calculated under the stated illustrative cost assumptions rather than observed retailer accounting outcomes.}
\label{fig:rq4}
\end{figure}
\FloatBarrier

\subsection{RQ5: Inventory Service-Level Trade-off}

RQ5 evaluates how estimated safety stock changes with the target service level. Table~\ref{tab:rq5} and Figure~\ref{fig:rq5} report the six service-level scenarios from 70\% to 95\%.
Using the observed LightGBM forecast-error dispersion and a seven-day lead time, estimated safety stock increases from 12.0 units at 70\% service to 37.8 units at 95\% service.
The intermediate estimates are 15.5, 19.3, 23.8, and 29.4 units at service levels of 75\%, 80\%, 85\%, and 90\%, respectively.
The relationship is monotonic because higher target service levels require larger buffers against forecast uncertainty.
The increase is not linear, reflecting the behaviour of the normal quantile in the safety-stock formulation.
The associated reorder-point simulation produces an illustrative reorder point of approximately 148.4 units under the 95\% service scenario.
Because the seven-day lead time is assumed rather than observed from a retailer replenishment system, these quantities are scenario estimates.
They should therefore not be interpreted as production inventory recommendations.

\begin{table}[!t]
\centering
\caption{Estimated safety-stock requirement by target service level using the observed M5 LightGBM forecast-error dispersion and a seven-day lead-time assumption.}
\label{tab:rq5}
\begin{tabular}{lc}
\toprule
\textbf{Target service level (\%)} & \textbf{Safety stock (units)} \\
\midrule
70 & 12.0 \\
75 & 15.5 \\
80 & 19.3 \\
85 & 23.8 \\
90 & 29.4 \\
95 & 37.8 \\
\bottomrule
\end{tabular}
\end{table}

\begin{figure}[!t]
\centering
\includegraphics[width=\linewidth]{Figures/rq5_service_level_safety_stock.pdf}
\caption{Estimated safety-stock requirement across target service levels from 70\% to 95\%. The estimates use the observed M5 LightGBM forecast-error dispersion and a seven-day lead-time assumption and are therefore scenario values.}
\label{fig:rq5}
\end{figure}
\FloatBarrier

\subsection{RQ6: Forecast Uncertainty and Stockout Risk}

RQ6 divides the 14,000 test observations into three equal-sized uncertainty tertiles using absolute forecast error. Table~\ref{tab:rq6} and Figure~\ref{fig:rq6} report the tier-specific error and safety-stock results.
The low-, medium-, and high-uncertainty groups contain 4,667, 4,666, and 4,667 observations, respectively.
Mean absolute errors are 1.036, 3.850, and 11.866 units across the three groups.
Using the corresponding error standard deviations, a 95\% service target, and a seven-day lead time produces safety-stock estimates of 2.74, 4.36, and 36.05 units.
The high-uncertainty group therefore requires a substantially larger illustrative buffer than the other groups.
The operational risk zones contain 8,541 low-risk, 4,534 medium-risk, and 925 high-risk observations under the notebook thresholds.
The relationship between absolute forecast error and stockout risk is mechanically positive because stockout risk is defined using the under-forecast component of demand error.
The result is therefore interpreted as a risk diagnostic rather than evidence of a causal relationship.

\begin{table}[!t]
\centering
\caption{Forecast uncertainty and error-based safety-stock estimates by uncertainty tertile. Safety stock uses a 95\% service target and seven-day lead time.}
\label{tab:rq6}
\begin{tabular}{lrrrr}
\toprule
\textbf{Tier} & \textbf{$n$} & \textbf{Mean abs. error} & \textbf{Std. dev.} & \textbf{Safety stock} \\
\midrule
Low & 4,667 & 1.036 & 0.631 & 2.74 \\
Medium & 4,666 & 3.850 & 1.001 & 4.36 \\
High & 4,667 & 11.866 & 8.283 & 36.05 \\
\bottomrule
\end{tabular}
\end{table}

\begin{figure}[!t]
\centering
\includegraphics[width=\linewidth]{Figures/rq6_uncertainty_safety_stock.pdf}
\caption{Forecast-error distribution and uncertainty-aware inventory diagnostics across the 14,000 M5 test observations. The high-error tail is associated with substantially larger estimated safety-stock requirements.}
\label{fig:rq6}
\end{figure}
\FloatBarrier

\subsection{RQ7: Robustness Across Demand Regimes}

RQ7 evaluates whether forecast performance remains stable across low-, medium-, and high-demand regimes. Table~\ref{tab:rq7} and Figure~\ref{fig:rq7} report the regime-level MAE results, while the table also records the separate DataCo robustness check.
The 500 selected series are divided using the training-period 28-day rolling-demand baseline.
LightGBM MAE increases from 3.595 in the low-demand regime to 5.155 in the medium-demand regime and 8.023 in the high-demand regime.
The high-demand regime therefore has more than twice the MAE of the low-demand regime.
The result shows that residual forecasting error is not distributed uniformly across the selected portfolio.
Higher-volume item--store combinations deserve particular attention when uncertainty-aware inventory policies are designed.
The rolling-MAE and two-dimensional density diagnostics support the same pattern without being interpreted as evidence of long-term deployment stability.
The conceptual scalability analysis indicates that the modular pipeline can be repeated for more forecasting entities, but no measured production-scale runtime benchmark is claimed. The separate DataCo external validation produces MAE 50.32 and RMSE 61.25 on its own target scale; these values are reported in Table~\ref{tab:rq7} as a robustness check rather than as a direct comparison with M5.

\begin{table}[!t]
\centering
\caption{LightGBM forecast accuracy by demand regime on the M5 test observations and the separate DataCo external-validation result. The DataCo metrics are reported on their own target scale and are not numerically compared with M5.}
\label{tab:rq7}
\begin{tabular}{llrr}
\toprule
\textbf{Analysis} & \textbf{Group / Dataset} & \textbf{MAE} & \textbf{RMSE} \\
\midrule
Demand regime & Low demand & 3.595 & -- \\
Demand regime & Medium demand & 5.155 & -- \\
Demand regime & High demand & 8.023 & -- \\
External validation & DataCo SMART SUPPLY CHAIN & 50.32 & 61.25 \\
\bottomrule
\end{tabular}
\end{table}

\begin{figure}[!t]
\centering
\includegraphics[width=\linewidth]{Figures/rq7_demand_regime.pdf}
\caption{LightGBM MAE across low-, medium-, and high-demand regimes. Forecast error increases with demand intensity, with the high-demand regime showing the largest MAE.}
\label{fig:rq7}
\end{figure}
\FloatBarrier

\subsection{Comparison with Existing Techniques}

The reported contemporary studies use different datasets, targets, forecast horizons, and evaluation protocols, so their numerical accuracy values are contextual rather than head-to-head comparisons \cite{Long2025ProbRetail,Ahmed2026RetailARIMA,Chintapanti2025PrGB,Punia2026HumanML}.
A supply-chain BiLSTM/NARX study reports MAE 1.73, RMSE 2.03, and $R^2=0.94$ on its own dataset \cite{Deepak2025RFLinear}.
An optimised XGBoost study reports $R^2=0.985$ for sustainable supply-chain performance prediction \cite{Nsir2025XGBoost}.
Probabilistic gradient-boosted approaches have been evaluated across M5, Favorita, and proprietary retail data using scaled probabilistic losses rather than the point-error metrics used here \cite{Long2025ProbRetail}.
M5-based work has also demonstrated that forecast accuracy and inventory performance can be related without being identical \cite{Theodorou2025Inventory}.
The present study does not claim state-of-the-art forecasting accuracy from these comparisons \cite{Abolghasemi2025LLMForecasting}.
Its contribution is instead the explicit connection of a reproducible forecasting pipeline to financial, inventory, uncertainty, demand-regime, and external-validation analyses using the same primary test observations.
This positioning is consistent with the study's objective of decision-oriented transparency rather than leaderboard optimisation.

\begin{table}[!t]
\centering
\caption{Reported accuracy of selected contemporary demand-forecasting studies alongside the present study. Values are contextual because datasets, targets, horizons, and evaluation protocols differ.}
\label{tab:comparison}
\scriptsize
\begin{tabular}{p{2.2cm} p{2.7cm} p{2.5cm} p{4.0cm}}
\toprule
\textbf{Study} & \textbf{Dataset scope} & \textbf{Best method} & \textbf{Reported result} \\
\midrule
Deepak et al. \cite{Deepak2025RFLinear} & Supply-chain demand & Uni-regression BiLSTM/NARX & MAE 1.73, RMSE 2.03, $R^2=0.94$ \\
Nsir et al. \cite{Nsir2025XGBoost} & Industry 4.0 supply chain & Optimised XGBoost & $R^2=0.985$ \\
Long et al. \cite{Long2025ProbRetail} & M5, Favorita, e-commerce & Probabilistic gradient-boosted trees & Dataset-dependent scaled pinball loss \\
Theodorou et al. \cite{Theodorou2025Inventory} & M5 competition & Multiple methods & Forecast--inventory relationship \\
\textbf{Present study} & \textbf{M5 top-500 item--store series} & \textbf{XGBoost} & \textbf{MAE 5.582, RMSE 8.690, $R^2=0.705$} \\
\bottomrule
\end{tabular}
\end{table}

% ============================================================
\section{Discussion}
% ============================================================

\subsection{RQ1: Feature Configuration}
The RQ1 results show that combining historical, pricing, and calendar information produced the lowest MAE among the four tested configurations.
The improvement from 5.842 to 5.608 is modest in absolute terms, while the calendar-only configuration deteriorates to 6.691.
This indicates that additional variables do not automatically improve a forecasting pipeline.
The result is more consistent with a feature-quality interpretation in which recent demand history provides the strongest signal and other sources add only limited incremental information.
The finding is relevant to structured retail forecasting because it supports a transparent feature-engineering approach rather than a simple increase in feature count.
It also provides a clear empirical basis for retaining the full feature configuration for the subsequent model comparison.
The conclusion is limited to the selected high-volume M5 series and the fixed chronological horizon used in this study.

\subsection{RQ2: Engineered Feature Contribution}
The RQ2 importance diagnostic shows that the 7-day rolling mean accounts for 79.9\% of the reported Random Forest importance, with the two rolling means together accounting for approximately 85.5\%.
This concentration indicates that recent demand history dominates the predictive signal in the selected feature matrix.
Price and calendar variables remain present but contribute much smaller shares.
The result should not be interpreted as causal evidence because lagged and rolling predictors are correlated.
Instead, the importance ranking is useful as a model diagnostic that identifies which information the fitted Random Forest relies on most strongly.
This provides a practical basis for monitoring the stability of temporal features when the pipeline is extended to other retail settings.
The finding also supports the decision to keep the model architecture relatively transparent rather than relying on a larger black-box feature space.

\subsection{RQ3: Core Model Comparison}
XGBoost provides the strongest row-level performance among the three core models, but its advantage over Random Forest is narrow.
The MAE difference is only 0.026 demand units, while both ensemble methods outperform Linear Regression on MAE and RMSE.
LightGBM produces a very similar result and was therefore suitable for the downstream analyses.
The supporting ARIMA result is not directly comparable because it is fitted after aggregation to a daily series.
The evidence therefore supports XGBoost as the strongest candidate for the specific feature matrix and horizon rather than as universally superior to other forecasting families.
This distinction is important because model rankings can change with demand structure, aggregation level, horizon, and feature availability.
The result is consequently best viewed as an empirical ranking under controlled conditions.

\subsection{RQ4: Financial Translation}
The RQ4 results demonstrate how the same held-out forecasts can be translated into scenario-level financial quantities.
Fulfilled-demand revenue is approximately \$335,984.96, while stockout-cost and holding-cost exposure are approximately \$6,020.06 and \$3,273.91.
The stockout-to-holding ratio is driven by the realised errors and the illustrative 10\% and 5\% cost factors.
It therefore cannot be treated as a general retail cost relationship.
The contribution of this analysis is the explicit connection between forecast output and decision-oriented financial quantities using the same test observations.
This makes the operational meaning of the forecast easier to inspect than a model score alone.
At the same time, retailer-specific costs would be required before these quantities could support a production financial decision.

\subsection{RQ5: Service-Level Trade-off}
The RQ5 results show a clear increase in estimated safety stock as the target service level rises.
The estimated requirement increases from 12.0 units at 70\% service to 37.8 units at 95\% service, with intermediate values increasing monotonically.
The non-linear increase reflects the normal-quantile component of the safety-stock formulation.
The result makes the service-level and inventory-investment trade-off explicit within the same forecasting framework.
The seven-day lead time is assumed, so the values should be interpreted as scenario estimates rather than observed retailer policy parameters.
Nevertheless, the pattern shows why a forecast can be operationally useful even when its accuracy is not perfect, provided its uncertainty is translated into a transparent inventory assumption.
The analysis therefore complements the predictive comparison with a decision-oriented view of forecast value.

\subsection{RQ6: Forecast Uncertainty and Stockout Risk}
The RQ6 results show that forecast uncertainty is highly uneven across the M5 test observations.
The high-uncertainty tertile has MAE 11.866 and estimated safety stock 36.05 units, compared with 1.036 and 2.74 units in the low-uncertainty tertile.
This supports differentiated inventory buffers rather than a uniform allowance across all observations.
The stockout-risk relationship is mechanically positive because stockout risk is defined from the under-forecast component of demand error.
It is therefore a diagnostic relationship rather than evidence of a causal effect.
The practical implication is that forecast-risk diagnostics can help identify observations for which inventory protection should be reviewed more closely.
The result also reinforces the importance of separating empirical risk diagnostics from prescriptive inventory policy.

\subsection{RQ7: Demand Regimes and External Robustness}
The RQ7 results show that MAE increases from 3.595 in the low-demand regime to 8.023 in the high-demand regime.
This indicates that forecast error is not uniform across the selected portfolio and that higher-demand series deserve particular attention when uncertainty-aware policies are considered.
The separate DataCo validation produces MAE 50.32 and RMSE 61.25 on a different target scale.
That result supports cautious transferability of the temporal feature-engineering strategy but does not establish cross-dataset metric equivalence or universal generalisation.
Taken together, the regime and external-validation results indicate that robustness should be examined both within the primary portfolio and across a structurally different dataset.
The conceptual scalability analysis further suggests that the modular pipeline can be repeated for more forecasting entities, although production-scale runtime has not been measured.
This keeps the interpretation aligned with the evidence actually generated by the notebook.

\subsection{Managerial Implications}
For retail planners, the findings support keeping recent demand history central to short-horizon forecasting while validating the incremental value of price and calendar variables.
The core model comparison indicates that a standard tree-based pipeline can provide strong predictive performance without requiring a deep-learning architecture.
The financial and safety-stock analyses show that forecasts should be accompanied by explicit assumptions about cost, lead time, and target service level before being converted into inventory decisions.
The uncertainty and demand-regime diagnostics can help prioritise observations where forecast risk is comparatively high.
The framework can therefore serve as a transparent decision-support layer rather than a fully prescriptive optimiser.
A production implementation would still require retailer-specific costs, lead-time distributions, replenishment constraints, substitution effects, and rolling-origin validation.
This interpretation is consistent with the study's objective of connecting predictive evidence to operational decision support without overstating the scope of the benchmark.

\subsection{Limitations}

The first limitation is the deliberate restriction to the 500 highest-volume item--store series.
The findings should not be automatically generalised to the full 30,490-series M5 hierarchy or to slow-moving and intermittent-demand products.
Second, the financial translation uses illustrative cost factors rather than retailer-specific procurement, shortage, holding, or lost-sales costs.
Consequently, the reported dollar quantities are scenario values rather than observed Walmart financial performance.
Third, the safety-stock calculations rely on an assumed seven-day lead time and simplified service-level formulation.
Actual replenishment lead-time distributions, constraints, substitution effects, and service definitions are not observed in M5.
Fourth, the DataCo experiment uses a different target scale and feature structure and is therefore not a controlled cross-dataset benchmark.
Finally, the scalability analysis is conceptual because the study does not report measured training time, memory consumption, or production deployment performance for very large numbers of forecasting entities.

\subsection{Future Directions}

Future research should extend the modelling sample beyond the top-500 high-volume series and evaluate whether the feature-engineering strategy remains effective for intermittent and low-volume demand.
A second direction is probabilistic forecasting, where prediction intervals or quantile forecasts can replace post-hoc uncertainty analysis and feed directly into inventory decisions.
A third direction is a richer inventory model using empirical lead-time distributions, retailer-specific cost parameters, service-level constraints, and replenishment policies.
Hierarchical forecast reconciliation could also be incorporated to test coherence across product, category, store, and geographical levels \cite{RIFT2026SupplyChain,Bergsma2025InventoryML}.
Measured computational benchmarks and rolling-origin evaluation across multiple forecast windows would provide stronger evidence for production-scale deployment and temporal robustness.
Multi-dataset validation could distinguish general feature-engineering effects from benchmark-specific behaviour.
Direct predictive-to-prescriptive optimisation is another potential extension because forecast-then-optimise pipelines can leave decision value on the table \cite{Bergsma2025InventoryML}.
These extensions would move the framework from transparent academic scenario analysis toward more complete operational decision support.

% ============================================================
\section{Conclusion}
% ============================================================

This study developed a reproducible machine-learning framework that connects retail demand forecasting with financial and inventory scenario analysis.
Using the M5 Forecasting dataset, the analysis selected the 500 highest-volume item--store series based only on training-period demand and evaluated a chronological 28-day hold-out containing 14,000 observations.
The full multi-source feature configuration achieved Random Forest MAE 5.608 compared with 5.842 for the lag-and-rolling baseline, while the 7-day rolling mean dominated the feature-importance diagnostic.
Among the three core models, XGBoost achieved the strongest row-level performance with MAE 5.582, RMSE 8.690, and $R^2=0.705$.
LightGBM produced a similar row-level result and was used for downstream decision-oriented analyses.
Under illustrative cost assumptions, the forecasts corresponded to approximately \$335,984.96 in fulfilled-demand revenue, \$6,020.06 in stockout-cost exposure, and \$3,273.91 in holding-cost exposure.
Estimated safety stock increased from 12.0 units at 70\% service to 37.8 units at 95\%, while uncertainty and demand-regime analyses showed that forecast risk is not uniform across the portfolio.
External validation on DataCo produced MAE 50.32 and RMSE 61.25 on its own scale, supporting cautious transferability of the feature-engineering approach.
Overall, the main contribution is not a claim of universal model superiority but a transparent framework that evaluates forecast accuracy and its operational implications from the same experimental evidence.

% ============================================================
\section*{Declarations}
% ============================================================

\subsection*{CRediT authorship contribution statement}
Mohamed Meraj Mansur Sharif: Conceptualization, Methodology, Software, Data curation, Formal analysis, Visualization, Writing -- original draft.\\
Raja Hashim Ali: Supervision, Validation, Writing -- review and editing.\\
Talha Ali Khan: Supervision, Validation, Writing -- review and editing.

\subsection*{Funding}
No external funding was reported for this study.

\subsection*{Declaration of Competing Interest}
The authors declare that they have no known competing financial or personal relationships that could have appeared to influence the work reported in this paper.

\subsection*{Data Availability}
The study uses publicly available benchmark and research datasets. The M5 Forecasting Accuracy dataset is publicly distributed for forecasting research through the M5 competition data repository. The DataCo SMART SUPPLY CHAIN dataset used for external validation is publicly available through Mendeley Data (DOI: 10.17632/8gx2fvg2k6.5).

\subsection*{Code Availability}
The experimental workflow was implemented in a Jupyter notebook executed in a Kaggle environment. The reproducibility notebook, figure-generation sources, and supporting LaTeX materials are included in the submission package accompanying this manuscript.

\subsection*{Declaration of generative AI and AI-assisted technologies in the manuscript preparation process}
During the preparation of this work, the authors used OpenAI ChatGPT for language editing and restructuring, assistance with selected passages, LaTeX formatting, and code and notebook debugging. The authors reviewed and edited the AI-assisted output and take full responsibility for the content of the published article.
The empirical research, data processing, model execution, reported results, and scientific interpretation were performed and verified by the authors against the project notebook and source files. Research and data visualizations in this manuscript were generated from the underlying data using reproducible analytical code; no generative-AI image tool was used to fabricate, alter, or replace research results. The graphical abstract was prepared separately using non-generative design tools.

\bibliographystyle{elsarticle-num}
\bibliography{Bibliography}

\end{document}
